\documentclass[11pt]{article}
\usepackage{amsmath,amssymb}
\usepackage{geometry}
\usepackage{enumitem}
\usepackage{multicol}
\geometry{margin=0.8in}
\setlength{\parindent}{0pt}
\setlength{\columnsep}{0.3in}

\begin{document}

\begin{center}
  {\Large\textbf{Complex Numbers: Polar Form, Powers, and Roots}}\\[3pt]
  {\large Euler's Formula and De Moivre's Formula --- 45 Minutes}
\end{center}

Name: \rule{2.7in}{0.4pt}\hfill Date: \rule{1.5in}{0.4pt}

\section*{Instructions}
Complete all 12 problems in 45 minutes without a calculator. You may use this
reference page. Show polar or exponential setup when computing powers and
roots. Give exact answers. Use rectangular form $a+bi$ unless another form is
requested. In final polar or exponential answers, use positive modulus and
angles in $[0,2\pi)$. When finding roots, list every distinct solution.

\section*{Reference Formulas}
\begin{multicols}{2}
\small
\textbf{Polar and exponential forms}

For $z=a+bi\ne0$,
\[
\begin{aligned}
  r=|z|&=\sqrt{a^2+b^2}>0,\\
  a&=r\cos\theta,\qquad b=r\sin\theta,\\
  z&=r(\cos\theta+i\sin\theta)=re^{i\theta}.
\end{aligned}
\]
Choose the angle using both coordinates and the correct quadrant.
Arguments differing by $2\pi k$, $k\in\mathbb Z$, represent the same direction.
Zero has modulus $0$ and no defined argument.

\medskip
\textbf{Euler's formula and Euler's identity}
\[
  e^{i\theta}=\cos\theta+i\sin\theta
\]
\[
  e^{i\pi}+1=0
\]
All angles are in radians. In particular, $e^{i(\theta+2\pi k)}=e^{i\theta}$.

\medskip
\textbf{Products, quotients, and conjugates}

For $z_j=r_je^{i\theta_j}$ with $r_j>0$,
\[
\begin{aligned}
  z_1z_2&=r_1r_2e^{i(\theta_1+\theta_2)},\\
  \frac{z_1}{z_2}&=\frac{r_1}{r_2}e^{i(\theta_1-\theta_2)},\\
  \overline z&=re^{-i\theta},\qquad
  \frac1z=\frac1r e^{-i\theta}.
\end{aligned}
\]

\columnbreak
\textbf{De Moivre's formula}

For $z=re^{i\theta}\ne0$ and integer $m$,
\[
\begin{aligned}
  z^m&=r^m e^{im\theta}\\
     &=r^m\bigl(\cos(m\theta)+i\sin(m\theta)\bigr).
\end{aligned}
\]
A power multiplies the argument by $m$ and raises the modulus to the
$m$th power. Negative powers require $z\ne0$.

\medskip
\textbf{All $n$th roots}

For $w=Re^{i\phi}\ne0$ and positive integer $n$, the equation $z^n=w$ has
exactly $n$ distinct solutions:
\[
  \boxed{z_k=R^{1/n}e^{\,i(\phi+2\pi k)/n}}
\]
\[
  k=0,1,\ldots,n-1.
\]
Equivalently,
\[
  z_k=R^{1/n}\left(\cos\frac{\phi+2\pi k}{n}
  +i\sin\frac{\phi+2\pi k}{n}\right).
\]
The roots lie on the circle $|z|=R^{1/n}$ and are equally spaced by an angle
of $2\pi/n$. For $n\ge3$, they are the vertices of a regular $n$-gon.

\medskip
The only distinct $n$th root of $0$ is $0$. To check any proposed root,
raise it to the $n$th power and recover $w$.
\end{multicols}

\newpage
\begin{center}
  {\large\textbf{Complex Numbers: Polar Form, Powers, and Roots --- Practice}}
\end{center}

\begin{enumerate}[leftmargin=*,itemsep=0pt]
  \item Write each number in trigonometric polar form
  $r(\cos\theta+i\sin\theta)$ and exponential form $re^{i\theta}$:
  \[
    z=-3\sqrt3+3i,\qquad w=-4-4i.
  \]
  \par\vspace{\stretch{1}}

  \item Evaluate exactly in rectangular form:
  \[
    e^{i\pi}+2e^{i\pi/2}-3e^{-i\pi/3}.
  \]
  \par\vspace{\stretch{1}}

  \item Simplify and give the result in both exponential and rectangular form:
  \[
    \frac{(2e^{5\pi i/6})(3e^{-\pi i/4})}{\sqrt3\,e^{\pi i/3}}.
  \]
  \par\vspace*{\stretch{1}}
\end{enumerate}

\newpage
\begin{center}
  {\large\textbf{Complex Numbers --- Continued}}
\end{center}

\begin{enumerate}[leftmargin=*,itemsep=0pt,resume]
  \item Express $\cos\theta$ and $\sin\theta$ using only $e^{i\theta}$,
  $e^{-i\theta}$, real constants, and $i$. Derive both formulas.
  \par\vspace{\stretch{1}}

  \item Evaluate in rectangular form:
  \[
    (-\sqrt3+i)^9.
  \]
  \par\vspace{\stretch{1}}

  \item Evaluate in rectangular form:
  \[
    \frac{(1+i)^8}{(\sqrt3-i)^6}.
  \]
  \par\vspace*{\stretch{1}}
\end{enumerate}

\newpage
\begin{center}
  {\large\textbf{Complex Numbers --- Continued}}
\end{center}

\begin{enumerate}[leftmargin=*,itemsep=0pt,resume]
  \item Let $z=2e^{i\pi/12}$. Find both quantities in rectangular form:
  \[
    z^{-6},\qquad \frac{z}{\overline z}.
  \]
  \par\vspace{\stretch{1}}

  \item Find all cube roots of $-27i$. Give every root in exponential
  and rectangular form.
  \par\vspace{\stretch{1}}

  \item Find all fourth roots of $-8+8\sqrt3\,i$. Give every root in
  exponential and rectangular form.
  \par\vspace*{\stretch{1}}
\end{enumerate}

\newpage
\begin{center}
  {\large\textbf{Complex Numbers --- Continued}}
\end{center}

\begin{enumerate}[leftmargin=*,itemsep=0pt,resume]
  \item Find all fifth roots of $32e^{5\pi i/6}$ in exponential form.
  State the radius of the circle containing the roots and their angular spacing.
  \par\vspace{\stretch{1}}

  \item Find all complex solutions in rectangular form:
  \[
    (z-1)^3=8i.
  \]
  \par\vspace{\stretch{1}}

  \item Find all complex solutions in exponential form:
  \[
    z^6+8z^3+64=0.
  \]
  \par\vspace*{\stretch{1}}
\end{enumerate}

\end{document}
